\chapter{Breast Pain}

\section*{Definition}
Breast pain is tenderness, aching, burning, or sharp discomfort in one or both breasts. It is often related to hormonal changes or chest-wall structures, but a new focal, persistent, or inflammatory pattern requires assessment.

\section*{When to See a Doctor}
Redness, warmth, fever, a rapidly enlarging lump, skin dimpling, nipple inversion, bloody discharge, or severe persistent focal pain. Arrange prompt evaluation for persistent, recurrent, unexplained, or function-limiting symptoms.

\section*{How It Develops}
Breast pain may reflect cyclic hormonal sensitivity, pregnancy or lactation, cysts, infection, trauma, chest-wall pain, medicines, or less commonly malignancy. Its relation to the menstrual cycle, focality, examination, and associated changes guide evaluation.

\section*{Causes}
\begin{itemize}
  \item Cyclical hormonal breast pain, pregnancy, breastfeeding, mastitis, cysts, trauma, medication effects, costochondritis, and rarely malignancy.
  \item Medication effects, environmental exposures, and underlying chronic disease may also contribute.
  \item Serious causes are less common but must be considered when symptoms are sudden, progressive, severe, or associated with systemic illness.
\end{itemize}

\section*{Related Symptoms}
\begin{itemize}
 \item Breast swelling, a lump, skin or nipple changes, or nipple discharge
 \item Fever, redness, warmth, or feeling unwell, especially during lactation
 \item Chest-wall tenderness, shortness of breath, or pain related to exertion
\end{itemize}

\section*{Medical Evaluation}
\begin{itemize}
  \item History addresses onset, duration, severity, triggers, injuries, exposures, medications, medical conditions, age, and pregnancy status when relevant.
  \item Examination includes the breasts and chest wall, with attention to focal tenderness, lumps, skin or nipple changes, and lymph nodes.
  \item Imaging is not usually needed for diffuse or cyclical pain without other suspicious findings. For focal, persistent, noncyclical pain, ultrasound is usually appropriate under age 30; diagnostic mammography or tomosynthesis and ultrasound may be appropriate from age 30 onward.
  \item Further tests are guided by the examination and imaging findings; a biopsy is not routine for pain alone.
\end{itemize}

\section*{Treatment Options}
Treatment may include a well-fitting supportive bra, simple analgesia when safe, and treatment of an identified infection, cyst, medicine effect, or chest-wall cause. Persistent focal pain or a mass requires clinical reassessment.

\section*{Self-Care and Home Management}
Avoid known triggers, protect the affected area, maintain appropriate hydration, and record timing and associated symptoms. Use only age- or pregnancy-appropriate medicines recommended by a clinician. Do not delay emergency care while trying home treatment.

\section*{References}
\begin{thebibliography}{99}
\bibitem{breast_pain:afp} Salzman B, Collins E, Hertz K. Common breast problems. \textit{Am Fam Physician}. 2012;86:343--349.
\bibitem{breast_pain:acog} American College of Obstetricians and Gynecologists. Benign Breast Disorders. ACOG Practice Bulletin No. 164. \textit{Obstet Gynecol}. 2016;128:e1--e16.
\bibitem{breast_pain:nice} National Institute for Health and Care Excellence. Suspected cancer: recognition and referral. NICE guideline NG12; 2015 (updated 2023).
\end{thebibliography}
